\section{工程与验证}

为了在真实工程中验证 Compound AI，讲者列出了多个 benchmark：Soang 游戏、Maze 路径搜索、以及纳米光学的 80\u00d780 网格设计。每一类都用 search trace 和 DPO 训练出的 agent（或 latent solver）与纯解法、纯生成模型做对比。

\subsection{Soang \& Maze 基准}

Soang 是一个只能推前的 box-pushing game，必须规划好顺序，否则一旦把 box 推入角落就回不去了。实验中，Solution-only baseline 需要约 1.75 亿参数、百万级训练样本才能偶尔成功；而 Search Former+DU Forer 在 15M 参数、10 万样本下即达到 80\% 正确率。
Maze benchmark 使用不同规模的 labyrinth：（1）sub 20 steps 的小型 maze；（2）延长到 40+ steps 的大图；（3）引入动态障碍物。Search trace 模型展示出比纯 LLM 更强的 generalization，尤其在分支繁多的 Maze 中，search trace 还提供有效的 candidate pruning。
\begin{knowledgebox}{工程结果摘要}
Soang 任务中的 trace dropout 提高了 robustness，Maze 任务则通过 consistency check（agent 自校的 judge）显著减少 hallucination。两个 benchmark 都表明：trace 模型用 5-10 倍更少的数据就能匹配甚至超越卷积的 solution-only baseline。
\end{knowledgebox}

\subsection{纳米光学 benchmark}

在光学设计任务中，目标是用 80\u00d780 网格的 binary pattern 控制每个像素的折射。训练 pipeline 需要：1）模拟频率响应；2）计算 beam splitter、wavelength multiplexing 的 loss；3）考虑制造时的 brush 限制。
讲者强调：只有把这些指标编入 loss，latent cost \(C\) 才能学习到既合法又高性能的设计。此外，将 solver 反馈回 \(C\) 的过程让模型可以自动从 failure trace 中学习，从而在不同 benchmark 之间迁移。

\subsection{本章小结}

这些 benchmark 既涵盖象棋式的 search，又涉及光学的连续设计，证明了 Compound AI 可以在多领域中一套模型串联多个功能模块，并通过简易的监督信号完成训练。

